\chapter[State of the art]{State of the art of the technology used or applied in this thesis}

Radio-map prediction sits at the intersection of classical wireless propagation
modelling, statistical channel analysis, and dense computer-vision regression.
This chapter surveys the relevant state of the art across four axes that jointly
motivate the design choices in Try~80: (i)~classical and empirical propagation
models, (ii)~deep learning for radio-map estimation, (iii)~physics-informed and
hybrid neural approaches, and (iv)~distributional and multi-target channel
modelling. A final section frames fair-comparison criteria and positions Try~80
within the literature.

%---------------------------------------------------------------------
\section{Channel knowledge maps as the thesis context}
\label{sec:soa_ckm_context}
%---------------------------------------------------------------------

The term channel knowledge map (CKM) refers to a site-specific database of
location-tagged channel knowledge, such as path loss, channel gain, angles, or
delay information, used to make future radio systems environment-aware
\cite{zengx2021ckm,ckmtutorial2024}. This framing is more precise than a
generic radio map: the map is not only an image-like coverage estimate, but a
queryable prior that can reduce channel-state acquisition overhead, guide
beamforming, support planning, and provide initial channel knowledge before
online measurements are available.

Recent CKM literature also clarifies why the thesis uses a large offline
ray-traced dataset rather than sparse field samples. Xu and Zeng
\cite{ckmdata2024} analyse how CKM construction error depends on spatial sample
density and on the number of neighbouring samples used for online prediction,
showing the explicit cost--accuracy trade-off of measurement-driven CKM
construction. For a thesis whose goal is to learn dense \(513\times513\)
targets under city holdout, this motivates the simulation-first route: it
provides full supervision everywhere, while the final model is still evaluated
as a surrogate for a ray-traced CKM rather than as a replacement for field
calibration.

The data ecosystem has also expanded since the first path-loss challenge.
RadioMap3DSeer \cite{dataset2212} and the ICASSP~2023 challenge
\cite{icassp2023challenge} provide the closest open benchmark family for
path-loss image prediction, while the directive-antenna dataset of Jaensch et
al.\ \cite{jaensch2024directiverme} and CKMImageNet
\cite{ckmimagenet2025} show the broader move toward open, location-specific
datasets with richer channel and environment representations. The CKM dataset
used in this thesis sits in the same data-driven tradition, but differs in its
continuous UAV height conditioning, real-city holdout split, and simultaneous
path-loss, delay-spread, and angular-spread targets.

%---------------------------------------------------------------------
\section{Classical propagation models and channel statistics}
\label{sec:soa_classical}
%---------------------------------------------------------------------

\subsection{Free-space and two-ray propagation}

The most fundamental path-loss term is the \emph{free-space path loss} (FSPL),
which depends only on carrier frequency~$f$ and Tx--Rx distance~$d$:
\begin{equation}
  \mathrm{PL}_{\mathrm{FSPL}} = 20\log_{10}(d) + 20\log_{10}(f)
    + 20\log_{10}\!\left(\frac{4\pi}{c}\right).
\end{equation}
FSPL is the irreducible floor in unobstructed conditions.  In aerial scenarios a
specular ground reflection produces constructive/destructive interference
periodically with distance.  The \emph{coherent two-ray} model captures this by
adding the direct and reflected fields before converting to dB
\cite{rappaport,wocc2021}:
\begin{equation}
  E_{\mathrm{total}} = E_{\mathrm{direct}}
    + \Gamma \cdot E_{\mathrm{reflected}} \cdot e^{j\Delta\phi},
\end{equation}
where $\Gamma$ is the reflection coefficient and $\Delta\phi$ is the phase
difference determined by the path-length difference.  A calibrated two-ray model
fitted to the CKM dataset (Try~78 of this thesis) achieves approximately
$1.75\,\mathrm{dB}$ RMSE on line-of-sight (LoS) pixels~--- the deterministic
radial structure explains most LoS variance without any neural network.

\subsection{Empirical and semi-empirical models}

Empirical models such as the Okumura--Hata and COST-231/Hata formulas extend
FSPL by incorporating frequency, antenna heights, and city-type correction terms
\cite{rappaport}.  While computationally cheap, they are calibrated on specific
frequency bands and morphology types, limiting accuracy when the geometry
departs from the fitting dataset.

The 3GPP Technical Report 38.901 \cite{tr38901} formalises stochastic channel
models for frequencies from 500~MHz to 100~GHz for urban macro (UMa), urban
micro (UMi), rural macro (RMa), and indoor scenarios.  It provides separate
path-loss polynomials for LoS and non-LoS (NLoS) conditions, breakpoint
distances, and shadow-fading standard deviations, which serve as reference
baselines throughout this thesis.

\subsection{UAV and air-to-ground channel models}

Classical terrestrial channel models assume the transmitter is at rooftop height
or below.  Unmanned aerial vehicle (UAV) base-station scenarios introduce a
third dimension: the UAV can fly from near-ground level to several hundred
metres.  Al-Hourani et al.~\cite{alhourani2014} showed that the LoS probability
increases monotonically with elevation angle and depends on city morphology type
(suburban, urban, dense urban, high-rise).  The ITU and 3GPP A2G channel model
family \cite{tr38901,khawaja_survey} parameterises this dependence through
logistic functions of the elevation angle
$\theta = \arctan\!\bigl((h_{\mathrm{tx}} - h_{\mathrm{rx}})/d_{2D}\bigr)$,
where $h_{\mathrm{tx}}$ and $h_{\mathrm{rx}}$ are transmit and receive heights
respectively and $d_{2D}$ is the horizontal distance.

Recent work by Saboor and Vinogradov \cite{saboor2025height} makes this
height dependence especially relevant for the present thesis: in simulated
26~GHz urban A2G links, LoS path loss remains close to a free-space exponent,
whereas NLoS path-loss exponents and shadow fading evolve with ABS height and
with urban layout. Vinogradov et al.\ \cite{vinogradov2026shadow} also propose
3D shadow projections as a fast way to generate spatially consistent A2G LoS
maps without full point-by-point ray tracing. These results support two design
choices used throughout this work: treating the LoS/NLoS mask as a physical
input rather than as a cosmetic label, and conditioning every dense prediction
on UAV height.

A key finding from the CKM dataset is that LoS path-loss residuals after
two-ray subtraction form almost perfect radial rings around the transmitter.
This confirms that buildings act primarily as a \emph{visibility mask} in LoS
conditions rather than as the main source of residual path-loss variance,
validating the prior-subtraction philosophy of Try~78.

\subsection{Delay spread and angular spread in standard channel models}

Beyond path loss, full channel characterisation requires second-order statistics.
The \emph{root-mean-square delay spread} $\tau_{\mathrm{rms}}$ measures the
time-domain dispersion of multipath arrivals and drives inter-symbol interference
in wideband systems.  The \emph{azimuth angular spread} $\sigma_{\mathrm{AS}}$
at the receiver quantifies the spatial diversity of arriving paths.

3GPP~TR~38.901 \cite{tr38901} models both quantities as log-normal random
variables conditioned on the LoS/NLoS state and distance.  The WINNER~II and
COST~2100 models adopt the same convention.  In practice, the distributions of
delay spread and angular spread are right-skewed with a heavy exponential tail
and often include a spike near zero representing the dominant LoS path.  Try~79
of this thesis empirically confirms this \emph{spike-plus-exponential} shape for
all topology classes in the CKM dataset, with skewness exceeding~5 and kurtosis
exceeding~40 in several morphology classes.  A Gaussian model cannot represent
this shape faithfully, motivating the mixture and distributional approaches of
Tries~77 and~80.

%---------------------------------------------------------------------
\section{Deep learning for radio-map estimation}
\label{sec:soa_dl}
%---------------------------------------------------------------------

\subsection{Image-to-image translation baselines}

The radio-map prediction problem can be cast as an image-to-image regression:
given an input topology map (and optionally a transmitter location channel), the
model predicts a path-loss map of the same spatial extent.  The pix2pix framework
\cite{isola2017pix2pix} established a general recipe for such tasks using a
conditional GAN with a U-Net generator and a PatchGAN discriminator.  Early
iterations of this thesis used this formulation.  In practice the adversarial
objective sharpens spatial texture but does not address the height-dependent
radial pattern or the heavy NLoS tails that dominate the error budget in the CKM
dataset.

\subsection{PMNet and the ICASSP outdoor challenge family}

PMNet \cite{pmnet2023} was the winner of the first ICASSP 2023 Outdoor Pathloss
Radio Map Prediction Challenge \cite{icassp2023challenge}.  It uses a deep
encoder-decoder with skip connections, trained and evaluated on the
RadioMap3DSeer dataset~\cite{dataset2212}.  RadioMap3DSeer consists of
procedurally generated city maps with transmitters fixed at rooftop height
(6--20~m) and building interior pixels masked to zero error.

The challenge metric is RMSE in normalised image units.  The RadioMap3DSeer
image mapping clips path gain to the range
$[\mathrm{PL}_{\mathrm{thr}}, \mathrm{PL}_{\mathrm{max}}] = [-111, -75]\,\mathrm{dB}$,
yielding a $36\,\mathrm{dB}$ dynamic window.  Converting to physical dB requires
multiplying normalised RMSE by~36.  After this correction, the true physical
errors are:

\begin{table}[ht]
\centering
\caption{Corrected physical RMSE for ICASSP~2023 outdoor challenge methods.
The normalised scores are taken from the challenge family and converted to dB
using the \SI{36}{dB} RadioMap3DSeer dynamic range described in the dataset
paper \cite{icassp2023challenge,dataset2212}.\protect\footnotemark}
\label{tab:icassp2023corrected}
\begin{tabular}{lcc}
\toprule
\textbf{Method} & \textbf{RMSE (norm.)} & \textbf{True RMSE (dB)} \\
\midrule
RMTransformer \cite{rmtransformer2025} & 0.0270 & 0.97 \\
PMNet \cite{pmnet2023} & 0.0383 & 1.38 \\
PPNet baseline \cite{icassp2023challenge} & 0.0507 & 1.83 \\
\bottomrule
\end{tabular}
\end{table}
\footnotetext{The RadioMap3DSeer dynamic range is defined by the path-gain
clipping window $[\mathrm{PL}_{\mathrm{thr}},\mathrm{PL}_{\mathrm{max}}] =
[-111,-75]\,\mathrm{dB}$ used to build the normalised target image, so the
physical span is $-75-(-111)=36\,\mathrm{dB}$. Multiplying the normalised RMSE
values reported in the challenge paper \cite{icassp2023challenge,dataset2212}
by this factor converts them to physical dB, which gives
$0.0383\times 36\approx 1.38$ for PMNet.}

These numbers are not directly comparable to the CKM evaluation for three
structural reasons: (i)~fixed Tx height eliminates the continuous UAV
height-conditioning challenge; (ii)~same-distribution train/test scenes allow
implicit generator-pattern memorisation; (iii)~the 36~dB window is much narrower
than CKM's effective 65--130~dB range, intrinsically compressing RMSE values.
The table is therefore useful as a scale reference for modern image-to-image
radio-map predictors, not as evidence that the Try~80 task is solved by those
architectures.  In particular, Try~80 must predict over unseen real-city
geometries, at native \(513\times513\) resolution, with a continuously varying
UAV altitude and with metrics computed only over valid ground receiver pixels.

A separate transfer test performed in this thesis (the ``weird tries'' baseline)
placed official PMNet checkpoints on CKM samples.  The best zero-shot result was
approximately $84\,\mathrm{dB}$ RMSE, with most models collapsing to a nearly
constant prediction.  After oracle linear calibration from test data the minimum
achievable RMSE was still approximately~$42\,\mathrm{dB}$.  This confirms that
PMNet's strong challenge score does not transfer to the structurally harder CKM
setting.

\subsection{Transformers, state-space models, and attention-based encoders}

The RMTransformer \cite{rmtransformer2025} replaced PMNet's convolutional
encoder with a multi-scale attention mechanism, improving the normalised RMSE on
RadioMap3DSeer to~0.0270 (true physical: $0.97\,\mathrm{dB}$).  The key
advantage is a larger receptive field for capturing long-range shadowing
correlations across wide urban streets.

RadioMamba and similar state-space sequence models applied to radio maps
\cite{wicopg2025,fmrme2026} offer linear complexity attention.  These
architectures are effective when global context is necessary, such as tracking
a shadowing corridor across an entire $513 \times 513$ map, but they require
substantially more data and compute than convolutional methods to reach the same
local detail quality.

\subsection{Foundation models for radio maps}

The most recent shift in the literature is towards \emph{wireless channel
foundation models} pre-trained on large heterogeneous datasets before
task-specific fine-tuning.  WiCo-PG \cite{wicopg2025} uses dual VQGANs and a
transformer with frequency-guided mixture-of-experts.  It exploits RGB drone
camera images as an auxiliary modality and achieves an NMSE of~0.012 on its
proprietary evaluation benchmark.  RadioLAM \cite{radiolam2025} targets fine-grained
3D radio map estimation with a large-model architecture at millisecond inference
latency.  FM-RME \cite{fmrme2026} frames radio map estimation as a foundation
modelling problem with self-supervised pre-training on simulated propagation
data.

These foundation approaches represent the long-term direction of the field but
require massive training corpora and compute budgets well beyond the scope of a
single CKM dataset.  They are cited here as the horizon against which this thesis
positions its own contributions.

%---------------------------------------------------------------------
\section{Physics-informed and hybrid neural approaches}
\label{sec:soa_physics}
%---------------------------------------------------------------------

\subsection{Geometry-assisted and diffraction-aware deep learning}

Purely data-driven models must implicitly learn diffraction geometry from the
topology map.  An alternative is to pre-compute a physics-based channel prior and
provide it as an additional input channel, reducing the learning problem to
residual estimation.

Shadow-projection methods \cite{vinogradov2026shadow} occupy an intermediate
point between stochastic LoS-probability models and full ray tracing. They use
3D building geometry to project blocked ground regions from the aerial
transmitter, producing a deterministic LoS map that is spatially consistent
along user trajectories. The CKM dataset already provides a ray-traced LoS
mask, so this thesis does not need to generate the mask itself; nevertheless,
the method is directly relevant because it explains why a binary visibility
channel can carry so much of the structure that the neural network otherwise
has to infer from the height map alone.

Knife-edge diffraction~\cite{geomDL2024} is a deterministic approximation
to single-edge scattering based on ITU-R~P.526 geometry.  For each pixel, the
first Fresnel zone obstruction height is computed from the 2D building
silhouette and converted to a diffraction gain.  This is a classical way of
injecting a physically meaningful NLoS cue into the input of a
geometry-assisted CNN.  Try~67 of this thesis implements this channel through
the \texttt{knife\_edge.py} module, feeding it into the PMHHNet encoder
alongside the topology map; the effect on the final Try~78--80 pipeline is
covered by the stronger calibrated NLoS prior rather than by a standalone
knife-edge ablation.

\subsection{Physics-informed neural networks}

The ReVeal architecture \cite{reveal2025} operates as a physics-informed neural
network (PINN) by incorporating the Helmholtz residual as an auxiliary loss term.
This penalises predictions that violate the wave equation, constraining the
network to produce physically smooth field maps and reducing corner-transition
artefacts.  ReVeal achieves~$1.95\,\mathrm{dB}$ RMSE in rural scenarios using
only 30 training points by leveraging the physical constraint as a strong
regulariser.  Tries~67--69 of this thesis implement a lightweight analogue:
a masked Laplacian penalty on the residual prediction, which penalises locally
non-smooth outputs without solving the full Helmholtz operator.

\subsection{Prior-residual decomposition}

A principled hybrid approach is to \emph{freeze} a calibrated analytic prior and
train the neural component to predict only the structured residual.  This design
reduces the effective learning problem from the full path-loss range to a
narrower residual distribution, improving data efficiency and interpretability.

Try~78 implements this philosophy for path loss: an enhanced two-ray model fitted
per height bin provides a prior map
$\hat{y}_{\mathrm{PL}}(\textbf{x}) = \mathrm{TwoRay}(d, h_{\mathrm{tx}})$,
and the neural component predicts the residual
$\epsilon = y - \hat{y}_{\mathrm{PL}}$.  Try~79 extends the idea to delay spread
and angular spread by using a ridge-regression model in log space with
TX-centred geometric features and local morphology statistics.  Try~80 unifies
all three targets under a single shared network anchored to both Try~78 and
Try~79 priors.

%---------------------------------------------------------------------
\section{Channel statistics prediction: delay spread and angular spread}
\label{sec:soa_stats}
%---------------------------------------------------------------------

\subsection{Empirical models for spread parameters}

Standard channel models characterise delay spread and angular spread as
log-normal distributions parameterised by distance and LoS/NLoS state.
TR~38.901 \cite{tr38901} provides polynomial fits for the log-normal mean and
standard deviation as functions of 2D distance and transmitter height across UMa,
UMi, and RMa scenarios.  These fits are limited to the morphology class and
frequency band of the underlying measurement campaigns.

Log-domain modelling is physically motivated: spread parameters are
multiplicative quantities in physical space, so their logarithm is more
Gaussian-like and additive-noise assumptions are better justified.  Try~79 adopts
this convention by working in $\log(1+y)$ space throughout, consistent with the
3GPP empirical modelling framework.

\subsection{Deep learning for wideband channel statistics}

Work on deep learning for delay and angular spread is sparser than for path
loss.  Most attempts either predict spread globally per sample (a single scalar)
rather than per pixel, or treat it as an ancillary output of a multitask
path-loss model.  Per-pixel dense prediction of delay spread and angular spread
over $513 \times 513$ maps is, to the best of the author's knowledge, not
addressed by published methods at the time of writing.

The main challenge is the spike-plus-exponential distribution shape of spread
targets.  A model trained only with mean-squared error (MSE) collapses to the
conditional mean, which lies in the interior of the spike and the exponential
tail and is therefore a poor representative of either regime.  Try~77 addresses
this with a Gaussian mixture model (GMM) distribution head, and Try~80
incorporates the Try~79 log-domain ridge prior as an anchor before training the
residual neural head.

\subsection{Distributional mismatch and the spike-plus-exponential problem}

The spike-plus-exponential shape~--- a Dirac-like concentration near zero
(the LoS-dominated case) plus a heavy exponential tail (NLoS multipath)~--- was
confirmed empirically on the CKM dataset.  Skewness exceeds~5 and excess kurtosis
exceeds~40 for both delay and angular spread in several topology classes.  A
single Gaussian or even a symmetric loss function does not fit this family.

Mixture models are the natural solution, but the required number of components
depends on what is being modelled.  The Try~76 histogram study selected the
best family for the raw target distributions, where path loss and the spread
targets often needed four or five components to match the full marginal shape.
Try~80 uses three components for a different object: the residual distribution
after the frozen Try~78/Try~79 priors have already explained the dominant LoS,
NLoS, delay-spread, and angular-spread structure.  This is a small design
change rather than a contradiction: the prior removes much of the target-side
multi-modality, so the residual head only needs a compact mixture to represent
under-correction, near-zero correction, and over-correction modes.  This
motivates the Stage-A distribution head in Try~76 (for path loss), Try~77 (for
spreads), and the lighter GMM-3 residual head carried into Try~80.

%---------------------------------------------------------------------
\section{Distribution-aware and multi-target modelling}
\label{sec:soa_distribution}
%---------------------------------------------------------------------

\subsection{Heteroscedastic uncertainty}

Kendall and Gal \cite{kendall2017uncertainties} established the framework for
heteroscedastic neural networks in computer vision: predict both a mean output
$\mu(x)$ and a log-variance $\log \sigma^2(x)$, optimising the negative
log-likelihood (NLL)
\begin{equation}
  \mathcal{L} = \frac{(\mu - y)^2}{\sigma^2} + \log \sigma^2.
\end{equation}
Pixels with high aleatoric uncertainty (deep NLoS shadows, building edges)
receive lower effective weight, while the network is penalised for
overestimating uncertainty everywhere.  Try~71 of this thesis applies this
loss to path-loss prediction, with $\sigma^2$ predicted as a second output
channel of PMHHNet.

\subsection{Gaussian mixture model heads}

Beyond heteroscedastic Gaussians, the GMM head allows multi-modal output
distributions.  Mixture-density networks have already appeared in wireless
channel modelling: Saleh et al.\ \cite{saleh2021probabilistic} use an MDN to
predict probabilistic mmWave path-loss distributions, Garcia Marti et al.\
\cite{garciamarti2020mixture} learn a stochastic Gaussian-mixture channel
model for physical-layer design, and Lee et al.\ \cite{lee2024timevarying}
derive conditional received-power PDFs with a mixture-density network.  These
works support the use of mixture outputs for wireless propagation, but they
operate at the level of scalar or channel-sample distributions rather than
dense radio-map generation.

Try~76 implements a two-stage architecture: Stage~A produces per-sample GMM
parameters $\{\pi_k, \mu_k, \sigma_k\}_{k=1}^{K}$ via a shallow convolutional
encoder with sinusoidal FiLM conditioning on UAV height; Stage~B is a small
U-Net that produces per-pixel component assignments $p_k(i,j)$ and
within-component deviations $z(i,j)$.  The final prediction is
\begin{equation}
  \hat{y}(i,j) = \sum_{k=1}^{K} p_k(i,j) \cdot \bigl(\mu_k + z(i,j) \cdot \sigma_k\bigr).
\end{equation}
This formulation avoids the mode-collapse failure of MSE regression while
maintaining a compact architecture.  The same two-stage GMM design is used in
Try~77 for delay and angular spread, with modified parameters for the
spike-plus-exponential prior structure.  The innovative aspect is therefore
not the isolated use of a Gaussian mixture, but its role as a global
distribution head that tells the dense decoder what value distribution the map
should contain before the model solves the harder pixel-placement problem.

\subsection{Stochastic diffusion models}

Conditional diffusion models applied to radio maps \cite{radiodiff2025} generate
physically plausible stochastic variations in deep shadow zones rather than
predicting a single blurry mean.  The iterative denoising procedure can be
conditioned on topology maps, transmitter locations, and height values.
RadioDiff reports $0.5$--$2\,\mathrm{dB}$ improvement over deterministic
regression baselines, primarily in NLoS-heavy regions where the posterior
distribution over path loss is genuinely multi-modal.  These models are
computationally expensive at inference time and are listed as a future-work
horizon for this thesis.

\subsection{Multitask and joint prediction}

Training a single shared backbone across multiple output targets can improve
sample efficiency and consistency between targets.  Path loss, delay spread, and
angular spread are physically correlated: a deep shadow corresponds to high
delay spread and potentially wide angular spread from diffracted paths.  A joint
model can therefore use the intermediate feature maps learned for path loss to
partially explain spread patterns, reducing the effective learning problem.

Try~80 instantiates this idea explicitly: one shared encoder processes the
topology map and UAV height, and three parallel decoder heads predict path-loss
residual, delay-spread residual, and angular-spread residual respectively, all
anchored to their corresponding physical priors.  The multi-task loss is a
weighted sum of per-target NLL terms, with weights tuned on the validation set.

%---------------------------------------------------------------------
\section{Training stabilisation and regularisation techniques}
\label{sec:soa_training}
%---------------------------------------------------------------------

\subsection{Stochastic weight averaging}

Stochastic weight averaging (SWA) \cite{izmailov2018swa} averages model weights
along the training trajectory at a low, approximately constant learning rate.
This finds wider optima in the loss landscape that generalise better to
unseen data.  The CKM city-holdout evaluation makes generalisation particularly
important since validation and test cities are never seen during training.
Tries~67--71 incorporate SWA starting from 60\,\% of the training budget.

\subsection{Exponential moving average weights}

Separate from SWA, exponential moving average (EMA) tracking of the model
weights at a decay rate of~0.9975 produces a smoother checkpoint suitable for
validation scoring.  EMA weights are used for validation and final test inference
but not for gradient updates.  The EMA-smoothed validation score is also used as
input to the ReduceLROnPlateau scheduler to avoid reacting to noisy single-epoch
validation spikes.

\subsection{Feature-wise linear modulation}

Feature-wise linear modulation (FiLM) \cite{wicopg2025,fmrme2026} conditions
the activations of a neural network on a scalar or vector input by applying
learned affine transforms $\gamma(c) \cdot \textbf{h} + \beta(c)$ to
intermediate feature maps.  Sinusoidal positional encoding of the UAV height
scalar (using log-spaced frequencies over the full 12--478~m range) is fed into
FiLM layers at multiple encoder depths.  This ensures that the UAV altitude
continuously modulates the activation statistics rather than being handled as a
discrete bin, which is critical when the effective propagation regime changes
substantially with height.

%---------------------------------------------------------------------
\section{Key benchmarks and fair-comparison protocol}
\label{sec:soa_benchmarks}
%---------------------------------------------------------------------

\subsection{Why headline numbers cannot be compared directly}

The most widespread mistake in radio-map literature reviews is comparing
headline RMSE values without verifying whether the underlying tasks are
equivalent.  The following evaluation axes are the minimum required for a
valid comparison:

\begin{enumerate}
  \item \textbf{Split:} city-holdout vs.\ in-distribution random split vs.\
        same-scene split.
  \item \textbf{Transmitter height:} fixed rooftop, discrete bins, or
        continuous UAV altitude.
  \item \textbf{Building pixels:} excluded from loss, excluded from metrics,
        set to zero error, or included.
  \item \textbf{Target quantity:} path loss (dB), path gain (dBm), received
        power, delay spread, angular spread.
  \item \textbf{Dynamic range and clipping:} determines the RMSE floor and
        ceiling independently of model quality.
  \item \textbf{Inference extras:} test-time augmentation (TTA), calibration,
        SWA, EMA weights; must be declared for both systems compared.
\end{enumerate}

This thesis uses a strict city-holdout split, continuous UAV height, ground-only
metrics (building pixels excluded from loss and metrics), native
$513 \times 513$ pixel maps, and no post-training calibration.  Table~\ref{tab:fairness}
summarises how published methods compare along these axes.

\begin{table}[ht]
\centering
\caption{Evaluation condition comparison across representative methods.}
\label{tab:fairness}
\small
\begin{tabularx}{\textwidth}{Xccccc}
\toprule
\textbf{Method} & \textbf{Unseen city} & \textbf{Cont.\ UAV height}
  & \textbf{$\ge$512\,px} & \textbf{Zero calib.} & \textbf{Comparable?} \\
\midrule
PMNet / ICASSP~2023 \cite{pmnet2023} & No & No & No & Yes & No \\
RMTransformer \cite{rmtransformer2025} & No & No & No & Yes & No \\
Gao et al.\ \cite{gao2026} & No & No & No & Yes & No \\
AIRMap \cite{airmap2025} & Partial & No & Yes & No & Partial \\
Tarhouni et al.\ \cite{tarhouni2025} & No & No & No & Yes & Physical ref. \\
PathFinder \cite{pathfinder2025} & Partial & No & No & Yes & Partial \\
\textbf{Try~80 (this thesis)} & \textbf{Yes} & \textbf{Yes} & \textbf{Yes}
  & \textbf{Yes} & \textbf{Reference} \\
\bottomrule
\end{tabularx}
\end{table}

\subsection{Tarhouni et al.\ --- measured suburban path-loss prediction}

Tarhouni et al.\ \cite{tarhouni2025} study machine-learning-based path-loss
prediction in the sub-\SI{6}{GHz} band using a suburban campus environment.  The
paper is valuable for this thesis because it sits closer to physical path-loss
prediction than purely image-domain RadioMap3DSeer challenge scores: the
reported errors are in dB and the setting is tied to a concrete environment
rather than only a procedural map generator.  However, it is not a strict
one-to-one comparison to CKM.  The setup is terrestrial/suburban, lower
frequency, and not a dense \(513\times513\) UAV CKM with per-pixel delay and
angular-spread targets.  This thesis therefore uses it as an outdoor path-loss
reference for the approximate difficulty of real-environment generalisation,
while keeping the final comparison qualitative.

The key lesson for Try~80 is not that the numerical RMSE should match
Tarhouni's setting, but that outdoor path-loss predictors evaluated outside
their training distribution tend to move back toward multi-dB errors even when
the same model family looks much better under in-distribution challenge
conditions.  This supports the decision to report Try~80 under a strict
city-holdout protocol.

\subsection{AIRMap --- adaptive digital-twin radio maps}

AIRMap \cite{airmap2025} is an adaptive U-Net autoencoder processing
$512 \times 512$ elevation maps for ultra-low-latency digital-twin
synchronisation.  It reports under $5\,\mathrm{dB}$ RMSE (path gain), but this
figure is obtained after lightweight field calibration from 20\,\% of the test
measurements.  The zero-shot baseline is also stated as under~$5\,\mathrm{dB}$
on Boston-only training, but the gap is plausibly explained by: (i)~calibrated
path gain vs.\ uncalibrated absolute path loss, (ii)~single-city training and
test vs.\ 66-city diversity, and (iii)~fixed terrestrial Tx vs.\ continuous UAV
height.

AIRMap is nevertheless an important citation for this thesis because it shows
that elevation-only deep surrogates can replace repeated ray tracing for
digital-twin style workloads.  Its reported inference time of about
\SI{4}{ms} on an NVIDIA L40S highlights the practical reason for the whole CKM
surrogate direction: the value of the model is not only lower RMSE, but the
ability to produce dense radio maps fast enough for planning loops.  Try~80 is
less mature than AIRMap as a deployable digital-twin system, but it attacks a
harder multi-target UAV setting and keeps physical priors visible instead of
using only a black-box elevation-to-path-gain network.

\subsection{Gao et al.\ --- multi-layer segmentation with corridor weighting}

Gao et al.\ \cite{gao2026} propose a ResNet-based outdoor path-loss predictor
with Tx depth maps, Rx depth maps, a distance map, and a weighting map that
emphasises the direct Tx--Rx corridor.  Using the ICASSP~2023 outdoor data,
Table~III of the paper reports their method at \SI{5.59}{dB} RMSE, compared
to \SI{8.09}{dB} for PPNet, \SI{9.56}{dB} for RPNet, and \SI{8.72}{dB} for a
twelve-layer Vision Transformer baseline; Table~II reports \SI{12.19}{dB}
RMSE on the ITU Challenge 2024 large-area subset, also ahead of the same
baselines, with roughly 60\% fewer FLOPs.  This is highly relevant to the
CKM work because it formalises a point
that appeared repeatedly in the internal tries: the physically important region
is not the whole topology image uniformly, but the subset of geometry that lies
on or near the propagation corridor.

Try~66 through Try~70 implement only a weaker Tx-centred radial proxy of that
idea.  Try~78 then solves the LoS part more cleanly with a two-ray radial prior,
while Try~80 uses the calibrated prior as an input channel.  A full Gao-style
per-Rx corridor map remains future work because generating one for every
receiver pixel at \(513\times513\) resolution is substantially more expensive
than using the current Tx-centred distance and obstruction features.

\subsection{PathFinder --- multi-transmitter domain adaptation}

PathFinder \cite{pathfinder2025} tackles distribution shift via disentangled
feature encoding and mask-guided low-rank attention, reporting the lowest MSE
of~$0.1068$ on unseen rural datasets.  The explicit domain-generalisation focus
aligns with the CKM city-holdout spirit, though PathFinder's setup involves
multi-transmitter inputs and a different normalisation convention.  Techniques
from PathFinder such as transmitter-oriented mixup augmentation are listed as
medium-priority future work for this thesis.

This citation matters because it separates two issues that are often mixed in
radio-map papers: prediction accuracy under the training distribution and
robustness under distribution shift.  PathFinder explicitly argues that many
methods over-focus on single-transmitter, in-distribution settings and
under-test multi-transmitter or shifted building-density regimes.  Try~80 is
not a multi-transmitter model, but its city-holdout split is motivated by the
same concern: a CKM surrogate is only useful if it works on new urban layouts,
not only on shuffled samples from the same city.

\subsection{ICASSP 2025 indoor challenge}

The ICASSP 2025 Indoor Pathloss Challenge \cite{icassp2025indoor} presents a
structurally different problem~--- dense wall-penetration NLoS effects in bounded
interior spaces~--- but its public results are a useful warning against
over-interpreting sub-\SI{2}{dB} outdoor challenge numbers.  The official result
page reports that the winning SIP2Net submission obtains a weighted-average RMSE
of \SI{9.411}{dB}; IPP-Net reaches \SI{9.501}{dB}; TerRaIn reaches
\SI{10.325}{dB}; and TransPathNet reaches \SI{10.397}{dB}
\cite{indoor2025results,sip2net2025,ippnet2025,transPathNet2025}.  These are
not outdoor UAV results, but they show that when NLoS physics becomes harder
and less directly observable from the input map, modern neural architectures
again hit a roughly \SIrange{9}{10}{dB} error band.

The comparison to Try~66 should therefore be read as diagnostic rather than
competitive.  Try~66's \SI{9.41}{dB} path-loss result is not claimed to beat
indoor methods; it is used to argue that the earlier CKM direct-regression
family had reached the same type of NLoS-dominated wall.  Try~78 and Try~80 were
developed precisely because improving that wall required a different
formulation: strong priors, regime separation, and residual distribution
modelling.

%---------------------------------------------------------------------
\section{Emerging techniques for closing the NLoS gap}
\label{sec:soa_emerging}
%---------------------------------------------------------------------

The $9$--$10\,\mathrm{dB}$ barrier common to both indoor and outdoor NLoS is not
a dataset scaling artefact; it reflects the fundamental difficulty of inferring
volumetric wave propagation from a 2D building height map.  Closing this gap
requires either richer inputs, stronger physical constraints, or stochastic
modelling of the aleatoric uncertainty.  Table~\ref{tab:emerging} summarises the
techniques with the highest estimated impact for CKM-style evaluation.

\begin{table}[ht]
\centering
\caption{Emerging techniques that could be applied on top of the current
pipeline, with implementation complexity as a rough guide. No numeric gain is
claimed here; each technique would need its own CKM ablation under the
city-holdout protocol to be quantified.}
\label{tab:emerging}
\begin{tabularx}{\textwidth}{clXc}
\toprule
\textbf{\#} & \textbf{Technique} & \textbf{Key reference} & \textbf{Complexity} \\
\midrule
1 & Knife-edge diffraction prior channel & \cite{geomDL2024} & Medium \\
2 & PINN Helmholtz residual loss & ReVeal \cite{reveal2025} & High \\
3 & Dual LoS/NLoS residual heads & Tries~69, 71 & Low \\
4 & Test-time adaptation on target city & RadioPiT \cite{radiopit2025} & Low--Med \\
5 & Diffusion-based residual refinement & RadioDiff \cite{radiodiff2025} & High \\
6 & Foundation-model pre-training & FM-RME \cite{fmrme2026} & Very high \\
\bottomrule
\end{tabularx}
\end{table}

Items~1--3 already exist in partial form inside the PMHHNet tries. Items~4--6
are roadmap items requiring additional data, compute, or infrastructure beyond
the scope of a single-GPU training setup.

%---------------------------------------------------------------------
\section{Position of Try 80 in the literature}
\label{sec:soa_position}
%---------------------------------------------------------------------

\subsection{What is novel about Try 80}

Try~80 occupies a distinct and innovative position in the radio-map literature
through the combination of four features that no published method located in
the thesis literature review addresses simultaneously:

\begin{itemize}
  \item \textbf{Distribution-first dense prediction.}  Prior MDN and GMM
        channel models predict scalar or low-dimensional propagation
        distributions \cite{saleh2021probabilistic,garciamarti2020mixture,lee2024timevarying}.
        Try~80 transfers that principle to dense radio maps: the model first
        estimates a compact sample-level distribution and then uses the
        decoder to place the corresponding modes spatially.  This directly
        follows the empirical observation that the predicted-vs-ground-truth
        histograms revealed the failure before the per-pixel maps alone did.

  \item \textbf{Three simultaneous dense targets.}  Try~80 jointly predicts
        per-pixel path loss, delay spread, and angular spread over
        $513 \times 513$ maps with a single shared backbone.  Published radio-map
        work addresses either path loss or spread statistics, but not all three
        in a pixel-dense city-holdout setting.

  \item \textbf{Frozen physics anchors.}  Two independently calibrated
        physics-based priors~--- the enhanced two-ray model (Try~78) for path
        loss and the log-space ridge regression (Try~79) for both spreads~--- are
        pre-computed and frozen.  The neural network predicts residuals with
        respect to these anchors rather than raw targets.  This reduces the
        effective learning range, improves interpretability, and allows the
        analytical prior to handle the bulk of the LoS structure without
        requiring the neural component to re-learn it.

  \item \textbf{Strict city-holdout evaluation with continuous UAV height.}
        All 66 CKM cities appear without repetition in train, validation, and
        test partitions.  The UAV transmitter height is a continuous scalar in
        $[12, 478]\,\mathrm{m}$ conditioned via sinusoidal FiLM, not a discrete
        bin.  This combination is not used in any published benchmark at the
        time of writing.
\end{itemize}

\subsection{What Try 80 does not claim}

Try~80 does not claim to replace full ray-tracing in arbitrary scenarios.  The
CKM dataset is a ray-traced simulation, and the thesis target is to reproduce
its outputs as accurately as possible under a strict generalisation protocol.
Physical effects that are not represented in the 2D topology inputs~--- such as
material-dependent penetration loss, antenna radiation patterns, or atmospheric
effects~--- are outside the prediction scope.

The thesis also does not claim parity with RadioMap3DSeer challenge numbers.
As discussed in Section~\ref{sec:soa_benchmarks}, those numbers apply to a
structurally easier evaluation contract (fixed height, same-distribution cities,
narrower dynamic range) and cannot be compared to CKM city-holdout RMSE without
systematic correction of all five axes listed in that section.

\subsection{Summary of position}

Table~\ref{tab:position} provides a compact summary of how Try~80 compares to
the most relevant published work across the evaluation dimensions introduced in
this chapter.

\begin{table}[ht]
\centering
\caption{Positioning of Try~80 against the most relevant published methods.}
\label{tab:position}
\scriptsize
\renewcommand{\arraystretch}{1.12}
\begin{tabularx}{\textwidth}{>{\raggedright\arraybackslash}p{0.28\textwidth}
                              >{\raggedright\arraybackslash}p{0.24\textwidth}
                              >{\centering\arraybackslash}p{0.11\textwidth}
                              >{\raggedright\arraybackslash}X}
\toprule
\textbf{Benchmark / method} & \textbf{Reported path-loss RMSE} & \textbf{City holdout} & \textbf{Interpretation} \\
\midrule
ICASSP~2023 best \cite{rmtransformer2025} & 0.97\,dB & No & Different target scaling and benchmark contract. \\
AIRMap \cite{airmap2025} (calibrated) & $<5$\,dB after calibration from 20\% of test measurements & Partial & Different protocol and path-gain target. \\
Gao et al.\ \cite{gao2026} (multi-layer) & 5.59\,dB on ICASSP~2023 data & No & Relevant corridor-weighting baseline, different protocol. \\
Tarhouni et al.\ \cite{tarhouni2025} & multi-dB, sub-\SI{6}{GHz} & No & Physical outdoor reference, not dense UAV CKM. \\
Try~66 (this thesis, single expert) & 9.41\,dB & Yes & Internal direct-regression baseline. \\
Try~78 non-DL prior & 1.92\,dB overall; 1.75\,dB LoS & Yes & Frozen path-loss prior used by the final pipeline. \\
\textbf{Try~80 (this thesis)} & \textbf{1.77\,dB validation snapshot; final test pending} & \textbf{Yes} & \textbf{This work; diagnostic until frozen test evaluation.} \\
\bottomrule
\end{tabularx}
\end{table}
\normalsize
The Try~80 number in Table~\ref{tab:position} is intentionally labelled as a
validation snapshot.  The final comparable thesis claim remains the frozen
Try~78 path-loss prior until the Try~80 checkpoint is selected and evaluated on
the held-out test split.  The prior-anchored design is expected to allow the
neural residual head to converge faster and to a narrower error floor than the
direct-regression family of Tries~66--71, particularly in LoS-dominated
morphology classes where the physics anchor already achieves sub-2\,dB error.
